\chapter{Turning candidate wavefronts into an engineering budget}\label{ch:budget}
\section{State the acceptance condition before assigning tolerances}
An engineering budget begins with an explicit list of requirements.
For pattern $p$, examples are
\[
|\mathrm{PS}_{y,p}|\leq b_{y,p},\quad
|\mathrm{CD}_{x,p}-\mathrm{CD}_{x,p}^{\rm target}|
\leq b_{{\rm CD},p},\quad
\mathrm{PVB}_{x,p}\leq b_{{\rm PVB},p}.
\]
The conditions should identify the field point, process range,
measurement definition, and whether compensation such as refocus or
global image registration is allowed.

The paper emphasizes a vertical-shift target near $\pm0.4\nm$ and keeps
other requested indicators in training-data ranges. That is a particular
candidate-generation experiment. A production budget would require
explicit physical limits for all relevant indicators, rather than
training coverage alone.

\section{A local constraint is a slab in coefficient space}
Let $\vct c_0$ be a verified candidate.
For indicator $i$, write its residual from the target as $r_i$ and its
gradient row as $\vct g_i^T$. Locally,
\[
r_i(\vct c_0+\delta\vct c)
\simeq r_i+\vct g_i^T\delta\vct c.
\]
The two-sided acceptance condition is
\begin{equation}
|r_i+\vct g_i^T\delta\vct c|\leq b_i.
\label{eq:slab}
\end{equation}
This is the region between two parallel hyperplanes.
Intersecting all such slabs gives a local feasible polytope.
An optical or mechanical constraint on coefficients further restricts it.
This geometry makes clear why tolerances are coupled.

\section{The largest local RMS ball around a candidate}
Suppose an orthonormal coefficient norm represents the retained RMS.
For any $\norm{\delta\vct c}\leq R$, Cauchy--Schwarz gives
\[
|\vct g_i^T\delta\vct c|
\leq\norm{\vct g_i}R.
\]
A sufficient local guarantee is
$|r_i|+\norm{\vct g_i}R\leq b_i$ for every $i$.
Hence a radius satisfying
\begin{equation}
R\leq\min_i\frac{b_i-|r_i|}{\norm{\vct g_i}}
\label{eq:ballbudget}
\end{equation}
fits inside the linearized constraints, omitting rows with zero gradient
unless their residual already violates the limit.

This is a bound on \emph{perturbations around a chosen candidate}.
It differs from the candidate's own RMS and from a maximum admissible
absolute RMS anywhere in the feasible set.
A system can have an appreciable baseline wavefront and still have
very little remaining tolerance to perturbation.

\section{Independent box limits require a joint worst-case test}
Suppose each coefficient perturbation is independently allowed within
$|\delta c_j|\leq a_j$. Then
\begin{equation}
\max_{|\delta c_j|\leq a_j}
|\vct g_i^T\delta\vct c|
=\sum_j|g_{ij}|a_j.
\label{eq:boxsupport}
\end{equation}
The upper bound follows from the triangle inequality.
It is attained by choosing the sign of each $\delta c_j$ to align
with $g_{ij}$.
A safe linear box therefore requires
\begin{equation}
|r_i|+\sum_j|g_{ij}|a_j\leq b_i
\quad\hbox{for every required metric }i.
\label{eq:boxbudget}
\end{equation}
Checking one coefficient at a time while fixing the others cannot
establish this joint guarantee.

\worked{The error in treating marginal ranges as a budget}
Suppose a required shift is proportional to $c_1+c_2$ and must obey
$|c_1+c_2|\leq0.4$ in chosen units.
The two candidates $(1,-1)$ and $(-1,1)$ both give zero shift.
Their observed coordinate ranges are $[-1,1]$ for both coefficients.
Combining the two maximum values gives $(1,1)$, whose shift is 2
and fails badly.

\figteach{feasible_budget}{.78}{A correlated feasible region and two
feasible candidates. The dashed box contains all independent mixtures of
the candidates' coordinate ranges. Its upper-right corner is not feasible.
This is why the min/max bars in the paper's Fig.~9(a) are not independent
guaranteed tolerances. The additional diagonal bound in this illustration
only makes the plotted region finite.}

\section{Ellipsoidal budgets and covariance}
For an ellipsoidal perturbation region
$\delta\vct c^TQ\delta\vct c\leq1$ with positive definite $Q$,
set $\vct z=Q^{1/2}\delta\vct c$. Then
\[
\vct g_i^T\delta\vct c
=(Q^{-1/2}\vct g_i)^T\vct z.
\]
Maximizing over $\norm{\vct z}\leq1$ gives
\begin{equation}
\max|\vct g_i^T\delta\vct c|
=\sqrt{\vct g_i^TQ^{-1}\vct g_i}.
\label{eq:ellipsoid}
\end{equation}
This allows tolerances to be wider along insensitive combinations and
tighter along sensitive combinations.

For random perturbations with covariance $C_c$,
linear propagation gives
\begin{equation}
C_y=JC_cJ^T,\qquad
\sigma_{y_i}^2=\vct g_i^TC_c\vct g_i.
\label{eq:covprop}
\end{equation}
The covariance formula is a statistical statement, not a deterministic
guarantee. If Gaussian assumptions are justified, a multiple of
$\sigma_{y_i}$ can be related to a tail probability.
For non-Gaussian tails or systematic errors, that interpretation may fail.

\section{A mode-specific placement calculation}
For our unit-RMS vertical tilt,
$W=c_yZ_1^{-1}=2c_yv$.
Equation \eqref{eq:tilt} gives
\begin{equation}
|\mathrm{PS}_y|=\frac{2|c_y|}{\NA}.
\end{equation}
At $\NA=0.33$, a pure-tilt requirement
$|\mathrm{PS}_y|\leq0.4\nm$ implies
\[
|c_y|\leq0.066\nm
=4.88889\,\mathrm{m}\lambda
\quad(\lambda=13.5\nm).
\]
A 20 milliwave tilt would shift the image by about 1.63636 nm.
This is an exact translation result within the stated scalar
shift-invariant pupil convention. It is not a numerical bound on the
paper's $Z_3$ because its normalization and axis convention are not given,
and compensating modes or registration could change the relevant allowance.

\section{From wavefront coefficients to physical hardware}
Let $\vct q$ collect physical design or adjustment variables:
mirror decentres, tilts, surface modes, or actuator commands.
Linearizing the optical system gives
\begin{equation}
\delta\vct c=B\,\delta\vct q,\qquad
\delta\vct y\simeq JB\,\delta\vct q.
\label{eq:hardware}
\end{equation}
$B$ is a physical influence matrix. Its column space contains the
wavefront changes actually attainable to first order.
A candidate coefficient vector outside that attainable set cannot be
implemented by the available controls, even if its simulated image is good.

A practical adjustment problem might minimize
\[
\norm{Q_y^{1/2}[F(\vct c_0+B\delta\vct q)-\vct y_*]}^2
+\alpha\norm{L_q\delta\vct q}^2
\]
subject to displacement, surface-slope, stress, alignment, or actuator
bounds. The paper does not supply $B$, so it does not solve this
hardware-realization step.

\worked{An illustrative mirror-height allocation}
Assume six equal, statistically independent mirror contributions, all
with local incidence angle $6^\circ$, negligible coating contribution,
and total permitted RMS OPD $0.27\nm$.
Using Eq.~\eqref{eq:mirror},
\[
\sigma_{W,\rm total}^2
=6(2\sigma_h\cos6^\circ)^2.
\]
Thus
\[
\sigma_h=\frac{0.27}{2\cos6^\circ\sqrt6}
=0.05542\nm.
\]
This is about 55 pm per surface under the stated assumptions.
Six perfectly correlated equal wavefront contributions would add in
amplitude and give a different allocation,
$0.27/[12\cos6^\circ]\simeq0.02262\nm$.
Neither number is a recovered tolerance of the projector in the paper.
They demonstrate how geometry and correlation enter a surface budget.

\section{Nonlinear robustness and process variation}
The Jacobian approximations are local. If a bound $M_i$ on the
relevant Hessian norm is known throughout a perturbation region, Taylor's
theorem gives a remainder estimate
\[
|R_i|\leq\frac12M_i\norm{\delta\vct c}^2.
\]
A conservative ball test can then include
$|r_i|+\norm{\vct g_i}R+\tfrac12M_iR^2\leq b_i$.
Without such a bound, a local linear result needs nonlinear checking
over the proposed tolerance region.

Likewise, checking focus at $-15,0,+15\nm$ does not mathematically
prove every intermediate focus passes. The extremum of a nonlinear
CD or PS curve can lie between samples.
PVB envelopes involve maxima and minima, and can become nonsmooth
when the process condition defining an extreme changes.

Useful engineering checks include nominal cases, candidate corners,
random perturbations from an explicit distribution, and targeted searches
for the worst violation. A finite test is evidence about the tested
region and sampler; it is not automatically a continuous worst-case proof.

\section{Finite success counts and probability}
As a teaching calculation, suppose $N$ cases are independently drawn
from a fixed operational distribution and every case passes.
If the true failure probability is $p$, the probability of seeing no
failures is $(1-p)^N$.
Setting this to 0.05 gives the usual one-sided 95 percent upper limit
\[
p_{\rm upper}=1-0.05^{1/N}.
\]
For $N=118$, this is about 2.51 percent.
Even under those strong sampling assumptions, zero observed failures
would not prove zero failure probability.
The paper does not establish that its selected 118 targets are
independent random samples from a manufacturing distribution, so this
calculation is not a statistical confidence claim about its results.

\section{Optical determinism does not remove photon and resist noise}
The paper's reported metrics are deterministic simulation indicators.
EUV photon energy is $E_\gamma=hc/\lambda\simeq91.84$ eV at 13.5 nm.
For a deliberately illustrative dose of $20\,\mathrm{mJ/cm^2}
=200\,\mathrm{J/m^2}$ incident on a $14\nm\times14\nm$ area,
\[
N_\gamma\approx
\frac{200(14\times10^{-9})^2}{hc/(13.5\times10^{-9})}
\simeq2.66\times10^3.
\]
This is an incident photon estimate for uniform exposure, not the actual
absorbed count in a printed contact image.

For ideal Poisson counting, $\Var N=\E N$ and relative count noise is
$1/\sqrt N$, about 1.94 percent for this example.
Absorption, secondary-electron transport, resist chemistry, and
nonuniform intensity add further structure.
The IRDS discussion identifies stochastic defects and pattern control
as major lithographic concerns \cite{irds}.
A wavefront that passes deterministic CD and PS criteria does not
thereby establish an acceptable missing-contact probability.

\section{A defensible use of the paper's method}
Use the inverse network as a candidate generator.
Verify each candidate with the forward model and all required patterns.
Check coefficient-domain coverage, physical realizability, and robustness
to manufacturing and process uncertainty. Then express the final budget
as an explicit joint set, distribution, or perturbation allowance with
documented conditions.

This preserves the paper's useful idea: imaging performance can guide
wavefront choices more directly than a single RMS number.
It also supplies the additional steps needed to turn promising examples
into an engineering tolerance statement.

\chapter{A reproducible laboratory for the key ideas}\label{ch:lab}
\section{What the supplied calculation does}
The script \texttt{tools/reproduce\_examples.py} generates this book's
original figures and the files \texttt{numerical\_checks.txt} and
\texttt{teaching\_model\_results.json}.
It verifies disk-mode orthogonality and several analytic identities,
then calculates a simple periodic contact image.
It contains no trained version of the authors' network and does not
claim to reproduce the commercial reflective-mask simulation.

The dependencies are NumPy, SciPy, and Matplotlib.
The PDF and figures already included with this book can be read without
running Python.

\section{Declared optical and numerical assumptions}
The contact example uses:
\begin{itemize}
\item wavelength 13.5 nm, image NA 0.33, and an unobscured unit pupil;
\item a scalar binary thin mask containing one centred
$14\nm\times14\nm$ opening per $38\nm\times38\nm$ period;
\item four idealized source sectors, radii 0.6--0.8, centred at
$45^\circ,135^\circ,225^\circ,315^\circ$;
\item 60 equally weighted source samples: three radii uniformly spaced
in squared radius and five midpoint angular samples per sector;
\item exact analytic rectangle Fourier coefficients from
Eq.~\eqref{eq:rectcoeff}, retaining every order that can enter the pupil
for the selected source support;
\item pupil phase $2\pi W/\lambda$ and paraxial defocus phase
$-\pi z\NA^2\rho^2/\lambda$;
\item a fixed threshold chosen so that the no-aberration, zero-focus,
unit-dose reference has a 14 nm CD on its central cuts.
\end{itemize}
Equal spacing in squared radius approximates uniform annular area
weighting because $\rho\dd\rho=\tfrac12\dd(\rho^2)$.
These choices make the example reproducible; they are not inferred
hidden settings from the paper.

\section{How the calculation is organized}
For each source point:
\begin{enumerate}
\item Compute normalized order locations
$\vct u_{mn}=\vct u_s+(m/p_x,n/p_y)/f_c$.
\item Reject orders with $|\vct u_{mn}|>1$.
\item Multiply each retained mask coefficient by its aberration and
defocus phase.
\item Sum the complex plane-wave orders along an image cut.
\item Square the field magnitude.
\end{enumerate}
Average those intensities across source points.
Find left and right crossings of the fixed threshold by linear
interpolation on a dense spatial grid.
Then calculate CD and midpoint shift.
For pure tilt, orthogonal cuts pass through the analytically translated
feature centre; otherwise a fixed off-centre cut could falsely suggest
a width change even for a rigid translation.

\section{Same RMS, different images}
The three nonzero cases use one unit-RMS mode with
$c=0.27\nm=20\,\mathrm{m}\lambda$:
horizontal tilt, defocus, and cosine astigmatism.
All have the same piston-removed wavefront RMS in the declared disk basis.

\begin{table}[htbp]\centering\small
\caption{Computed teaching-model results, rounded for readability.
The exact arrays are saved with the source. These are not paper data.}
\begin{tabular}{@{}lrrrr@{}}
\toprule Case & CD$_x$ (nm) & PS$_x$ (nm)
& CD$_y$ (nm) & PS$_y$ (nm)\\ \midrule
Reference &14.00000&0&14.00000&0\\
Horizontal tilt &14.00000&$-1.63636$&14.00000&0\\
Defocus &13.94249&0&13.94249&0\\
Astigmatism &13.97132&0&13.97132&0\\
\bottomrule
\end{tabular}
\end{table}
Tilt agrees with the analytic displacement $-2c/\NA$ and preserves
the registered width. Defocus and astigmatism change the image profile
by different amounts.
In this particularly symmetric in-focus example, the astigmatism case
has equal central-cut CDs in the two axes. That is not a universal
statement about astigmatism; different focus, source, or mask conditions
can break this equality.

\figteach{contact_profiles}{1.0}{Original computed contact profiles.
Left: equal-RMS phase modes have different effects on a fixed-threshold
image. Right: a close-up of the right edge reveals the small change in
threshold crossing through focus. The $+15$ nm and $-15$ nm curves
coincide by symmetry in this example.
All curves use the same reference threshold rather than retuning dose
to conceal profile changes.}

\section{Focus and dose variation}
The demonstration samples focus at $-15,0,+15\nm$ and dose at
$0.95,1.00,1.05$. It collects left and right edge positions over all
nine combinations.
Both edge-band widths are approximately
\[
B_L=B_R=0.974586\nm
\]
in this symmetric teaching example.
The equality follows from the model symmetry, not from a universal
PVB identity. Because the centre remains fixed here,
the full CD range is twice one edge-band width.

\figteach{focus_dose_bands}{1.0}{Left: teaching Bossung curves computed
over a wider focus interval for illustration. Right: left and right
edge positions for the nine stated focus-dose samples. Shaded regions
show their individual process bands.}

\section{Independent checks built into the script}
The script uses Gauss--Legendre integration in $\rho^2$ and uniform
angular quadrature to check the Gram matrix of all 28 real modes
through radial degree six. Its largest numerical departure from
the identity is about $5.1\times10^{-14}$ in the supplied run.
It separately checks the balanced-quartic variance, the analytic
tilt displacement, preservation of registered widths under tilt, and
simple inverse-ambiguity and correlated-budget counterexamples.

It also prints the wavelength conversions, defocus RMS, 2033-row
arithmetic, hypothetical full-grid size, architecture parameter count,
and mask-scale consistency estimates. These checks test the calculations
actually used in this book.

The source sample count is intentionally modest for transparent,
quick execution. There is no claim that these pedagogical images are
electromagnetically rigorous or numerically converged to a production
accuracy requirement. Exact identities are checked tightly; approximate
optical-model output is labelled according to its assumptions.

\section{Experiments you can perform after reading}
\begin{enumerate}
\item Change pitch while keeping the hole width fixed. Track how the
admitted order set changes and identify which changes in the image
follow from altered interference.
\item Change the sign of a pure tilt coefficient. Confirm that the
shift reverses and registered widths stay fixed.
\item Add a small focus offset to astigmatism. Examine whether horizontal
and vertical CDs remain equal.
\item Reduce edge slope by changing source or threshold, then apply the
same dose perturbation. Compare measured edge motion with
Eq.~\eqref{eq:nilsdose}.
\item Increase source and spatial sampling. Check which plotted values
stabilize and how long the calculation takes.
\item Replace a single-mode perturbation by two modes together. Compare
the combined metric change with the sum of separate changes to reveal
nonlinear interactions.
\end{enumerate}
These exercises teach the mechanism behind the paper's dataset.
They do not require pretending that the missing commercial simulation
or neural-network weights are available.
